% =============================================================================
% METHODOLOGY CHAPTER: OBSEA Multivariate Data Processing Pipeline
% PhD Thesis - Coastal Ocean Monitoring and AI-Based Time Series Reconstruction
% =============================================================================
\documentclass[12pt,a4paper]{article}

% --- Essential Packages ---
\usepackage[utf8]{inputenc}
\usepackage[english]{babel}
\usepackage{graphicx}
\usepackage{amsmath,amssymb}
\usepackage{booktabs}
\usepackage{hyperref}
\usepackage{float}
\usepackage{caption}
\usepackage{subcaption}
\usepackage[left=2.5cm,right=2.5cm,top=2.5cm,bottom=2.5cm]{geometry}
\usepackage{setspace}
\usepackage{enumitem}
\usepackage{xcolor}
\usepackage{algorithm}
\usepackage{algpseudocode}
\usepackage{tabularx}
\usepackage{multirow}
\usepackage{longtable}

% Hyperlink configuration
\hypersetup{
    colorlinks=true,
    linkcolor=blue!70!black,
    citecolor=blue!70!black,
    urlcolor=blue!70!black
}

\onehalfspacing

% Custom commands
\newcommand{\degC}{$^\circ$C}
\newcommand{\psu}{PSU}

% Better figure positioning
\renewcommand{\floatpagefraction}{0.8}
\renewcommand{\topfraction}{0.9}
\renewcommand{\bottomfraction}{0.9}
\renewcommand{\textfraction}{0.1}

\begin{document}

% =============================================================================
\section{Methodology: Multivariate Oceanographic Data Processing and AI-Based Gap Filling}
\label{sec:methodology}
% =============================================================================

\subsection{Overview}

This chapter presents a comprehensive methodology for the integration, quality control, and reconstruction of oceanographic time series data from the OBSEA underwater observatory (Vilanova i la Geltrú, NW Mediterranean Sea). The pipeline processes heterogeneous sensor data from 2009 to 2024, producing a gap-free, scientifically validated multivariate dataset at 30-minute temporal resolution.

The methodology follows a sequential, modular architecture comprising seven distinct phases:

\begin{enumerate}[noitemsep]
    \item \textbf{Data Acquisition}: Loading raw observations from multiple instruments.
    \item \textbf{Temporal Harmonization}: Resampling to a canonical 30-minute grid.
    \item \textbf{Instrumental Quality Control}: QARTOD-compliant data validation.
    \item \textbf{Physical Feature Engineering}: Derivation of oceanographic variables.
    \item \textbf{Gap Analysis and Classification}: Identification and categorization of missing data.
    \item \textbf{AI-Based Imputation}: Machine learning reconstruction of gaps.
    \item \textbf{Validation and Export}: Quality assessment and dataset generation.
\end{enumerate}

\begin{figure}[H]
    \centering
    \includegraphics[width=\textwidth]{figures/obsea_pipeline_diagram.png}
    \caption{High-level architecture of the OBSEA data processing pipeline. The workflow proceeds from raw sensor acquisition (left) through quality control and physical feature extraction to AI-based gap filling, culminating in a Digital Twin visualization platform (right).}
    \label{fig:pipeline_overview}
\end{figure}

\newpage
% =============================================================================
\subsection{Phase 1: Data Acquisition}
\label{sec:data_acquisition}
% =============================================================================

The OBSEA observatory provides data from five primary instrument categories, each with distinct sampling characteristics and measurement domains (Table~\ref{tab:data_sources}).

\begin{table}[H]
\centering
\caption{Summary of OBSEA data sources, instruments, and measured variables. The dataset spans 2009--2024.}
\label{tab:data_sources}
\begin{tabularx}{\textwidth}{@{}llX@{}}
\toprule
\textbf{Instrument} & \textbf{Type} & \textbf{Primary Variables} \\ \midrule
CTD SBE37 & Oceanographic & Temperature (TEMP), Salinity (PSAL), Conductivity (CNDC), Pressure (PRES) \\
AWAC (Currents) & Hydrodynamic & Current speed (CSPD), Direction (CDIR), U/V components (UCUR, VCUR) \\
AWAC (Waves) & Hydrodynamic & Significant wave height ($H_s$), Peak period ($T_p$), Mean direction \\
Airmar & Meteorological & Wind speed (WSPD), Direction (WDIR), Air temperature (ATMP), Barometric pressure \\
CTVG Vantage Pro2 & Meteorological & Land-based reference: wind, humidity, solar radiation \\
\bottomrule
\end{tabularx}
\end{table}

% =============================================================================
\subsection{Phase 2: Temporal Harmonization}
\label{sec:harmonization}
% =============================================================================

Raw observations are resampled to a canonical 30-minute grid. The approach is \textbf{physically coherent}, applying different aggregation methods based on variable type:

\begin{itemize}[noitemsep]
    \item \textbf{Scalar variables} (temperature, salinity): Arithmetic mean.
    \item \textbf{Angular variables} (wind/wave direction): Circular mean (Equation~\ref{eq:circular_mean}).
    \item \textbf{Wave periods}: Median (robust to outliers).
\end{itemize}

\begin{equation}
    \bar{\theta} = \arctan2\left(\frac{1}{n}\sum_{i=1}^{n} \sin\theta_i,\, \frac{1}{n}\sum_{i=1}^{n} \cos\theta_i\right)
    \label{eq:circular_mean}
\end{equation}

This decomposition into sine and cosine components correctly handles the 0°/360° angular discontinuity.

% =============================================================================
\subsection{Phase 3: Instrumental Quality Control (QARTOD)}
\label{sec:quality_control}
% =============================================================================

Quality control follows the international QARTOD framework through four sequential automated tests:

\begin{enumerate}
    \item \textbf{Range Check}: $x_{min} \leq x_i \leq x_{max}$ (e.g., $10\degC \leq T \leq 28\degC$)
    
    \item \textbf{Spike Detection}: A value is flagged if:
    \begin{equation}
        |x_i - \bar{x}_{window}| > \lambda \cdot \sigma_{window}, \quad \lambda = 3
        \label{eq:spike}
    \end{equation}
    
    \item \textbf{Gradient Check}: Rate of change must satisfy:
    \begin{equation}
        \left|\frac{x_i - x_{i-1}}{\Delta t}\right| \leq \left(\frac{dx}{dt}\right)_{max}
    \end{equation}
    
    \item \textbf{Flatline Detection}: Consecutive identical values exceeding $n_{threshold} = 6$ (3 hours).
\end{enumerate}

% =============================================================================
\subsection{Phase 4: Physical Feature Engineering}
\label{sec:feature_engineering}
% =============================================================================

The pipeline derives physically meaningful variables using established oceanographic relationships:

\begin{table}[H]
\centering
\caption{Derived physical variables and their scientific basis.}
\label{tab:derived_vars}
\begin{tabularx}{\textwidth}{@{}lXl@{}}
\toprule
\textbf{Variable} & \textbf{Description} & \textbf{Standard} \\ \midrule
$\sigma_\theta$ & Potential density anomaly: $\rho(S_A, \Theta, 0) - 1000$ & TEOS-10 \\
$N^2$ & Brunt-Väisälä frequency (stratification proxy) & — \\
$\tau$ & Wind stress: $\rho_{air} \cdot C_d \cdot U_{10}^2$ & Bulk formula \\
Wave Energy & Power flux: $\frac{\rho g^2 H_s^2 T_p}{64\pi}$ [kW/m] & Linear theory \\
Anomaly & Deviation from climatological mean: $x' = x - \bar{x}(\text{doy}, \text{hour})$ & — \\
\bottomrule
\end{tabularx}
\end{table}

\textbf{Stationarity Analysis:} The Augmented Dickey-Fuller (ADF) test identifies non-stationary series ($p \geq 0.05$).

\textbf{STL Decomposition:} Seasonal-Trend decomposition using LOESS separates:
\begin{equation}
    Y_t = T_t + S_t + R_t
\end{equation}
where $T_t$ is the long-term trend (providing annual change rates), $S_t$ is the seasonal component, and $R_t$ is the residual.

\newpage
% =============================================================================
\subsection{Phase 5: Gap Analysis and Classification}
\label{sec:gap_analysis}
% =============================================================================

Missing data segments are systematically identified and categorized by temporal duration. This taxonomy (Table~\ref{tab:gap_taxonomy}) drives the \textbf{automatic selection of imputation methods}, ensuring that each gap is handled by the most appropriate algorithm.

\begin{table}[H]
\centering
\caption{Gap classification taxonomy and method selection rationale.}
\label{tab:gap_taxonomy}
\begin{tabularx}{\textwidth}{@{}llXX@{}}
\toprule
\textbf{Category} & \textbf{Duration} & \textbf{Selected Method} & \textbf{Rationale} \\ \midrule
Micro   & $< 1$ hour       & Linear interpolation  & Local continuity preserved; minimal physical change expected \\
Short   & 1--6 hours       & VARMA                 & Cross-variable covariance captures oceanographic relationships \\
Medium  & 6 hours--3 days  & Bi-LSTM               & Non-linear temporal dependencies; diurnal cycles captured \\
Long    & 3--30 days       & XGBoost (bidirectional) & Recursive ensemble exploits both past and future context \\
Extended & 30--60 days     & XGBoost + climatology & Seasonal information required; physics-guided constraints \\
Gigant  & $> 60$ days      & Climatological fill   & Beyond machine learning predictability horizon \\
\bottomrule
\end{tabularx}
\end{table}

\subsubsection{Method Selection Logic}

The pipeline implements an automated decision tree for imputation method selection:

\begin{algorithm}[H]
\caption{Automatic Method Selection Based on Gap Duration}
\begin{algorithmic}[1]
\Require Gap duration $d$, configuration flags
\If{$d < 1$ hour}
    \State Apply \textbf{Linear Interpolation}
\ElsIf{$d < 6$ hours \textbf{and} VARMA enabled}
    \State Apply \textbf{VARMA} with correlated predictors
\ElsIf{$d < 3$ days \textbf{and} Bi-LSTM enabled}
    \State Apply \textbf{Bi-LSTM with Attention}
\ElsIf{$d < 60$ days \textbf{and} XGBoost enabled}
    \State Apply \textbf{XGBoost Bidirectional Recursive}
\Else
    \State Apply \textbf{Climatological Mean} fill
\EndIf
\end{algorithmic}
\end{algorithm}

\begin{figure}[H]
    \centering
    \includegraphics[width=\textwidth]{figures/gaps_heatmap.png}
    \caption{Data availability heatmap across all variables and the 15-year observation period. Green indicates data presence; red indicates missing data. Vertical patterns reveal instrument failures affecting multiple variables simultaneously.}
    \label{fig:gap_heatmap}
\end{figure}

\begin{figure}[H]
    \centering
    \includegraphics[width=0.85\textwidth]{figures/gap_duration_histogram.png}
    \caption{Distribution of gap durations. Left: Gaps under 24 hours (linear scale). Right: All gaps (log scale) with category boundaries marked.}
    \label{fig:gap_histogram}
\end{figure}

\begin{figure}[H]
    \centering
    \includegraphics[width=\textwidth]{figures/timeseries_ctd_gaps.png}
    \caption{CTD time series (Temperature, Salinity, Pressure) with classified gaps overlaid as colored bands. Green = micro, yellow = short, orange = medium, red = long, purple = extended.}
    \label{fig:ctd_gaps}
\end{figure}

\newpage
% =============================================================================
\subsection{Phase 6: AI-Based Gap Filling Models (MTSI Framework)}
\label{sec:imputation}
% =============================================================================

The reconstruction of missing data follows the \textbf{Multivariate Time Series Imputation (MTSI)} framework. The core principle of MTSI in this study is the simultaneous exploitation of two orthogonal dynamics:
\begin{itemize}[noitemsep]
    \item \textbf{Longitudinal Dynamics:} The temporal inertia and autocorrelation of a single variable over time.
    \item \textbf{Cross-sectional Dynamics:} The instantaneous physical covariance between different variables (e.g., Temperature, Salinity, and Sound Velocity).
\end{itemize}

\subsubsection{Methodological Design Principles}
To ensure the scientific rigor of the MTSI process before deploying advanced algorithms, the pipeline adheres to strict methodological principles:

\begin{enumerate}
    \item \textbf{Missingness Masking:} The neural architectures incorporate a binary mask ($M \in \{0,1\}$) concatenated to the input features. This explicitly informs the model which inputs are genuine observations and which are zero-filled proxies, allowing the network to dynamically adjust its attention and trust.
    \item \textbf{Strict Chronological Splitting:} Unlike standard cross-validation, temporal order is strictly preserved. No random shuffling is applied during the Train/Validation/Test split to prevent future data leakage into the past.
    \item \textbf{Simulation-based Benchmarking:} Performance is evaluated by simulating artificial gaps (MCAR, MAR) with varying distributions over periods of known, high-quality data.
    \item \textbf{Physics-Informed Loss Functions:} Deep learning models are trained using custom loss functions that penalize not only the Mean Squared Error (MSE) but also unphysical gradients and values falling outside normalized climatological bounds.
\end{enumerate}

The pipeline integrates eleven different algorithms ranging from classical statistics to the 2024 state-of-the-art in deep learning, mapped to specific gap durations as defined in Section \ref{sec:gap_analysis}. Three of the most representative machine learning approaches are detailed below.

% -----------------------------------------------------------------------------
\subsubsection{Model 1: VARMA (Vector Autoregressive Moving Average)}

For short gaps where linear cross-correlations dominate, VARMA exploits multivariate structure:

\begin{equation}
    \mathbf{y}_t = \sum_{i=1}^{p} \Phi_i \mathbf{y}_{t-i} + \sum_{j=1}^{q} \Theta_j \boldsymbol{\epsilon}_{t-j} + \boldsymbol{\epsilon}_t
\end{equation}

where $\mathbf{y}_t$ is the observation vector and $\Phi_i$, $\Theta_j$ are coefficient matrices. Model order $(p,q)$ is selected via information criteria (AIC/BIC).

% -----------------------------------------------------------------------------
\subsubsection{Model 2: Bidirectional LSTM with Attention}

For medium and long gaps, a deep learning architecture captures non-linear dependencies.

\paragraph{Architecture.} The model consists of:
\begin{itemize}[noitemsep]
    \item Two stacked Bidirectional LSTM layers
    \item Attention mechanism for temporal weighting
    \item Fully connected output layer
\end{itemize}

\begin{table}[H]
\centering
\caption{Bi-LSTM model hyperparameters and their values.}
\label{tab:lstm_params}
\begin{tabular}{@{}llp{7cm}@{}}
\toprule
\textbf{Hyperparameter} & \textbf{Value} & \textbf{Description} \\ \midrule
\texttt{hidden\_size} & 128 & Number of LSTM units per direction (256 total for bidirectional) \\
\texttt{num\_layers} & 2 & Stacked LSTM layers for hierarchical feature extraction \\
\texttt{dropout} & 0.2 & Regularization between layers to prevent overfitting \\
\texttt{bidirectional} & True & Processes sequence forward and backward \\
\texttt{sequence\_length} & 96 & Input window: 96 steps $\times$ 30 min = 48 hours \\
\texttt{batch\_size} & 32 & Mini-batch size for stochastic gradient descent \\
\texttt{epochs} & 100 & Maximum training iterations \\
\texttt{learning\_rate} & 0.001 & Adam optimizer step size \\
\texttt{early\_stopping\_patience} & 10 & Stop if validation loss doesn't improve for 10 epochs \\
\texttt{train\_ratio} & 0.6 & 60\% of data for training \\
\texttt{val\_ratio} & 0.2 & 20\% for validation, 20\% for testing \\
\bottomrule
\end{tabular}
\end{table}

\paragraph{Attention Mechanism.} The model learns to weight the importance of each time step:

\begin{equation}
    e_t = \mathbf{v}^T \tanh(\mathbf{W}[\mathbf{h}_t; \mathbf{s}] + \mathbf{b}), \quad
    \alpha_t = \frac{\exp(e_t)}{\sum_{k=1}^{T} \exp(e_k)}
\end{equation}

where $\mathbf{h}_t$ are encoder hidden states and $\mathbf{s}$ is the decoder state. The context vector $\mathbf{c} = \sum_t \alpha_t \mathbf{h}_t$ focuses on relevant time steps.

\paragraph{Residual Learning.} The model is trained to predict the \textit{anomaly} rather than the raw signal:

\begin{equation}
    \hat{x}' = f_{LSTM}(\mathbf{X}), \quad \hat{x} = \hat{x}' + \bar{x}_{climatology}
\end{equation}

This allows the network to focus on non-trivial variability.

\paragraph{Blending Correction.} A linear bias ramp ensures smooth transitions at gap boundaries:

\begin{equation}
    \hat{x}_i^{corrected} = \hat{x}_i + \text{bias}_{start} \cdot \left(1 - \frac{i}{n}\right) + \text{bias}_{end} \cdot \frac{i}{n}
\end{equation}

where $\text{bias}_{start/end}$ are the prediction errors at the first and last known values.

% -----------------------------------------------------------------------------
\subsubsection{Model 3: XGBoost Bidirectional Recursive}

For production-scale imputation, XGBoost provides efficient gradient boosting with optimized recursive prediction.

\begin{table}[H]
\centering
\caption{XGBoost model hyperparameters and their values.}
\label{tab:xgb_params}
\begin{tabular}{@{}llp{7cm}@{}}
\toprule
\textbf{Hyperparameter} & \textbf{Value} & \textbf{Description} \\ \midrule
\texttt{n\_estimators} & 500 & Number of boosting rounds (trees) \\
\texttt{max\_depth} & 7 & Maximum tree depth; controls model complexity \\
\texttt{learning\_rate} & 0.05 & Shrinkage factor per iteration \\
\texttt{subsample} & 0.8 & Row sampling ratio for each tree \\
\texttt{colsample\_bytree} & 0.8 & Feature sampling ratio for each tree \\
\texttt{min\_child\_weight} & 3 & Minimum sum of instance weight in a child \\
\texttt{gamma} & 0.1 & Minimum loss reduction for split \\
\texttt{reg\_alpha} & 0.1 & L1 regularization (Lasso) \\
\texttt{reg\_lambda} & 1.0 & L2 regularization (Ridge) \\
\texttt{device} & cuda & GPU acceleration (falls back to CPU) \\
\texttt{tree\_method} & hist & Histogram-based algorithm (GPU-efficient) \\
\bottomrule
\end{tabular}
\end{table}

\paragraph{Feature Engineering.} The model uses engineered temporal features:

\begin{table}[H]
\centering
\caption{Temporal features used by XGBoost.}
\label{tab:xgb_features}
\begin{tabular}{@{}lp{9cm}@{}}
\toprule
\textbf{Feature Type} & \textbf{Values} \\ \midrule
Lag features & 2, 4, 12, 24, 48, 168 steps (1h to 1 week) \\
Rolling windows & 12, 48, 168 steps with mean and std \\
Cyclical features & $\sin/\cos$ encoding of hour, day-of-week, day-of-year \\
Time features & Hour, day, month, year, is\_weekend \\
\bottomrule
\end{tabular}
\end{table}

\paragraph{Recursive Prediction Algorithm.}

\begin{algorithm}[H]
\caption{XGBoost Recursive Gap Filling (Optimized $O(n)$)}
\begin{algorithmic}[1]
\Require Trained model $M$, time series with gap, feature matrix $\mathbf{F}$
\State Pre-compute static features (cyclical, exogenous) once
\State Identify gap positions $\{i_1, i_2, \ldots, i_n\}$
\For{each position $i_k$ in gap}
    \State Extract feature row $\mathbf{x}_{i_k}$ from $\mathbf{F}$
    \State Predict: $\hat{y}_{i_k} = M(\mathbf{x}_{i_k})$
    \State Update dynamic features (lags, rolling stats) with $\hat{y}_{i_k}$
\EndFor
\State \Return Filled time series
\end{algorithmic}
\end{algorithm}

\paragraph{Bidirectional Fusion.} Two models are trained—forward and backward. Final predictions are averaged:

\begin{equation}
    \hat{x}_i = \frac{\hat{x}_i^{forward} + \hat{x}_i^{backward}}{2}
\end{equation}

This cancels directional bias and improves accuracy in gap centers.

% -----------------------------------------------------------------------------
\subsubsection{Model 4: State-of-the-Art Deep Learning (MTSI Suite)}

To guarantee the highest scientific fidelity, especially for extended sensor failures, the pipeline has been upgraded to include an MTSI suite of foundational models:

\begin{itemize}
    \item \textbf{BRITS (Bidirectional Recurrent Imputation for Time Series):} Unlike standard RNNs, BRITS explicitly models the temporal decay -- the time elapsed since the last valid observation -- forcing the hidden state to regress towards the mean during long gaps.
    \item \textbf{SAITS (Self-Attention-based Imputation for Time Series):} A Transformer architecture that replaces recurrence with a joint-optimization of masked imputation and standard reconstruction tasks, capturing long-range dependencies efficiently.
    \item \textbf{ImputeFormer:} A 2024 model that introduces a \textit{Low-Rankness} inductive bias into the Transformer architecture, making it exceptionally robust for multivariate datasets where spatial (cross-variable) correlations dominate the missingness pattern.
    \item \textbf{MissForest:} A non-parametric, Random Forest-based imputer included as a rigorous machine learning baseline for tabular relationships.
\end{itemize}

\newpage
\subsubsection{Imputation Examples by Gap Category}

\begin{figure}[H]
    \centering
    \includegraphics[width=0.9\textwidth]{figures/example_gap_micro.png}
    \caption{\textbf{Micro gap ($<$ 1 hour):} Linear interpolation. The short duration allows simple connection of endpoints without introducing artifacts.}
    \label{fig:gap_micro}
\end{figure}

\begin{figure}[H]
    \centering
    \includegraphics[width=0.9\textwidth]{figures/example_gap_short.png}
    \caption{\textbf{Short gap (1--6 hours):} VARMA imputation. Cross-correlation with related variables (e.g., air temperature) improves reconstruction.}
    \label{fig:gap_short}
\end{figure}

\begin{figure}[H]
    \centering
    \includegraphics[width=0.9\textwidth]{figures/example_gap_medium.png}
    \caption{\textbf{Medium gap (6 hours--3 days):} Bi-LSTM with attention. The model captures diurnal variability and non-linear dynamics.}
    \label{fig:gap_medium}
\end{figure}

\begin{figure}[H]
    \centering
    \includegraphics[width=0.9\textwidth]{figures/example_gap_long.png}
    \caption{\textbf{Long gap (3--30 days):} XGBoost bidirectional recursive. The forward-backward ensemble provides robust predictions even for extended missing periods.}
    \label{fig:gap_long}
\end{figure}

\newpage
% =============================================================================
\subsection{Phase 7: Validation and Benchmarking}
\label{sec:validation}
% =============================================================================

Model performance is assessed through \textbf{block-masking cross-validation}:

\begin{enumerate}[noitemsep]
    \item Artificial gaps are created by masking known values.
    \item Each imputation method fills the masked segment.
    \item Predictions are compared against ground truth.
\end{enumerate}

\subsubsection{Evaluation Metrics}

\begin{align}
    \text{RMSE} &= \sqrt{\frac{1}{n}\sum_{i=1}^{n}(y_i - \hat{y}_i)^2} \label{eq:rmse} \\[0.5em]
    \text{MAE} &= \frac{1}{n}\sum_{i=1}^{n}|y_i - \hat{y}_i| \label{eq:mae} \\[0.5em]
    R^2 &= 1 - \frac{\sum_{i=1}^{n}(y_i - \hat{y}_i)^2}{\sum_{i=1}^{n}(y_i - \bar{y})^2} \label{eq:r2}
\end{align}

\subsubsection{MTSI Post-Processing}
Following the predictive phase, MTSI requires a post-processing stage to ensure the numerical outputs translate into valid oceanographic states:
\begin{itemize}[noitemsep]
    \item \textbf{Denormalization:} Inverse scaling of standard normal predictions back to their physical units.
    \item \textbf{Physical Bounding (Clipping):} Hard constraints restricting predictions within instrument limits and theoretical bounds (e.g., $Salinity > 0$, directional values modulo $360^{\circ}$).
    \item \textbf{Boundary Blending:} Application of a weighted cross-fade over the edge samples connecting the imputed gap to the genuine observation to prevent artificial, high-frequency steps in the continuous series.
\end{itemize}

\begin{figure}[H]
    \centering
    \includegraphics[width=\textwidth]{figures/benchmark_metrics_RMSE.png}
    \caption{Benchmark results: RMSE comparison across imputation methods and gap categories for Temperature reconstruction. Lower values indicate better performance. XGBoost and Bi-LSTM consistently outperform traditional methods for medium and long gaps.}
    \label{fig:benchmark}
\end{figure}

% =============================================================================
\subsection{Phase 8: Output Generation and Visualization}
\label{sec:output}
% =============================================================================

The final outputs include:

\begin{itemize}[noitemsep]
    \item \textbf{Reconstructed dataset}: \texttt{OBSEA\_multivariate\_30min\_interpolated.csv}
    \item \textbf{Imputation tracking}: Metadata indicating which method filled each gap
    \item \textbf{Statistical summaries}: Descriptive statistics and correlation matrices
    \item \textbf{Visualizations}: Time series plots, T-S diagrams, gap heatmaps
\end{itemize}

\begin{figure}[H]
    \centering
    \begin{subfigure}[b]{0.48\textwidth}
        \includegraphics[width=\textwidth]{figures/correlation_matrix.png}
        \caption{Variable correlation matrix.}
    \end{subfigure}
    \hfill
    \begin{subfigure}[b]{0.48\textwidth}
        \includegraphics[width=\textwidth]{figures/ts_diagram.png}
        \caption{T-S diagram with $\sigma_\theta$ isopycnals.}
    \end{subfigure}
    \caption{Diagnostic visualizations: (a) Pearson correlation between oceanographic and meteorological variables reveals physical relationships exploited by multivariate models. (b) Temperature-Salinity diagram characterizes water mass properties.}
    \label{fig:diagnostics}
\end{figure}

\begin{figure}[H]
    \centering
    \includegraphics[width=\textwidth]{figures/comparison_full_timeseries.png}
    \caption{Full 15-year time series comparison: Original observations (blue) versus AI-reconstructed continuous dataset (orange). Gaps are seamlessly filled while preserving physical variability patterns.}
    \label{fig:full_comparison}
\end{figure}

% =============================================================================
\subsection{Computational Environment}
\label{sec:environment}
% =============================================================================

\begin{table}[H]
\centering
\caption{Software environment and hardware specifications.}
\label{tab:environment}
\begin{tabular}{@{}ll@{}}
\toprule
\textbf{Component} & \textbf{Specification} \\ \midrule
Python version & 3.10 \\
Deep learning & PyTorch 2.1 with CUDA 12.1 \\
Gradient boosting & XGBoost 2.0 (GPU-enabled) \\
Oceanography & GSW library (TEOS-10) \\
Statistics & statsmodels, scipy \\
GPU & NVIDIA GeForce RTX 4060 Ti (8 GB) \\
\bottomrule
\end{tabular}
\end{table}

GPU acceleration reduces LSTM training time from days to hours and enables efficient XGBoost recursive prediction.

% =============================================================================
\subsection{Summary}
\label{sec:summary}
% =============================================================================

This methodology transforms heterogeneous, gap-riddled sensor data into a high-fidelity, continuous multivariate dataset suitable for climate analysis. The combination of rigorous quality control (QARTOD), physics-based feature engineering (TEOS-10), and state-of-the-art machine learning (Bi-LSTM, XGBoost) ensures both scientific validity and practical utility.

The resulting dataset spans 15 years at 30-minute resolution, providing an unprecedented resource for studying coastal ocean dynamics in the NW Mediterranean Sea.

% =============================================================================
\section*{References}
% =============================================================================

\begin{enumerate}[label={[\arabic*]}]
    \item IOC, SCOR and IAPSO. \textit{The International Thermodynamic Equation of Seawater – 2010} (TEOS-10). UNESCO, 2010.
    \item U.S. IOOS. \textit{Manual for Real-Time Quality Control of In-Situ Temperature and Salinity Data} (QARTOD). 2020.
    \item Chen, T. and Guestrin, C. \textit{XGBoost: A Scalable Tree Boosting System}. KDD'16, 2016.
    \item Hochreiter, S. and Schmidhuber, J. \textit{Long Short-Term Memory}. Neural Computation, 1997.
    \item Bahdanau, D., Cho, K., and Bengio, Y. \textit{Neural Machine Translation by Jointly Learning to Align and Translate}. ICLR, 2015.
\end{enumerate}

\end{document}
